\section{Finite-memory conjugacy and dynamical consequences}\label{sec:measurement}

Theorem~\ref{thm:two-sided-conjugacy} turns the sharp threshold into a dynamical statement. Above threshold, stabilization is not merely a bounded-degree sofic factor; it is a finite-memory conjugacy onto a finite-type image. This section records the resulting consequences for block presentations, path-space equivalence, and thermodynamic formalism.

\subsection{Block presentations above the threshold}\label{sec:measurement:sufficiency}

The theorem below is best viewed as a deterministic equivalence between raw and stabilized block presentations. The language of sufficiency is a corollary of this bijection, not the primary structural input.

Fix $m\ge 3$ and $n\ge m$. Write
\[
\mathcal Y_{m,n}:=\Phi_{m,n}\bigl(\{0,1\}^{n+m-1}\bigr)\subset X_m^n
\]
for the legal length-$n$ stabilized output blocks, and let
\[
\Psi_{m,n}:\mathcal Y_{m,n}\to\{0,1\}^{n+m-1}
\]
denote the inverse of $\Phi_{m,n}$ from Theorem~\ref{thm:block-invertibility}.

\begin{theorem}[Deterministic equivalence of block presentations above threshold]\label{thm:exact-sufficiency}
Fix $m\ge 3$ and $n\ge m$. Let $\{P_u:u\in\mathcal U\}$ be any family of probability laws on the raw block space $\{0,1\}^{n+m-1}$, and let
\[
Q_u:=(\Phi_{m,n})_*P_u
\]
be the induced family on $\mathcal Y_{m,n}$. Then:
\begin{enumerate}[label=(\roman*),leftmargin=*,itemsep=2pt]
  \item For every $u\in\mathcal U$, if $A^{(n)}\sim P_u$ and $Y^{(n)}:=\Phi_{m,n}(A^{(n)})$, then
  \[
  A^{(n)}=\Psi_{m,n}(Y^{(n)})\qquad P_u\text{-a.s.}
  \]
  Consequently,
  \[
  H_{P_u}(A^{(n)}\mid Y^{(n)})=0.
  \]
  \item For any random parameter $U$ with conditional law $A^{(n)}\mid\{U=u\}\sim P_u$,
  \[
  I(U;Y^{(n)})=I(U;A^{(n)}).
  \]
  \item The raw-block and stabilized-block experiments are deterministically equivalent. In particular, for every $u,v\in\mathcal U$,
  \[
  D_{\KL}(P_u\|P_v)=D_{\KL}(Q_u\|Q_v),
  \qquad
  D_{\TV}(P_u,P_v)=D_{\TV}(Q_u,Q_v),
  \]
  the Chernoff information agrees, and every Bayesian or Neyman--Pearson decision problem has the same optimal risk under the two experiments.
\end{enumerate}
\end{theorem}

\begin{proof}
Theorem~\ref{thm:block-invertibility} states that $\Phi_{m,n}$ is injective, so $\Psi_{m,n}$ is well-defined on $\mathcal Y_{m,n}$ and satisfies
\[
\Psi_{m,n}(\Phi_{m,n}(a))=a
\qquad\text{for every }a\in\{0,1\}^{n+m-1}.
\]
Hence $A^{(n)}=\Psi_{m,n}(Y^{(n)})$ almost surely under every $P_u$, proving (i).

For (ii), $Y^{(n)}$ is a deterministic function of $A^{(n)}$, so by data processing
\[
I(U;Y^{(n)})\le I(U;A^{(n)}).
\]
Conversely, $A^{(n)}$ is a deterministic function of $Y^{(n)}$ by (i), so the reverse inequality also holds.

For (iii), each $Q_u$ is the pushforward of $P_u$ under a bijection between $\{0,1\}^{n+m-1}$ and $\mathcal Y_{m,n}$. Therefore any divergence or decision criterion that is invariant under deterministic bijections is identical for the two experiments. The displayed identities for $D_{\KL}$ and $D_{\TV}$ are obtained by rewriting the defining sums on $\mathcal Y_{m,n}$ via the inverse map $\Psi_{m,n}$. Chernoff information is preserved for the same reason, and a decision rule on one experiment transports to the other by composition with $\Phi_{m,n}$ or $\Psi_{m,n}$ without changing its loss distribution.
\end{proof}

\subsection{Path-space equivalence and the localization of irreversibility}

The point of this subsection is to separate the fixed-window aliasing ledger from any claim about process-level loss. Above threshold the stabilized process is measurably isomorphic to the raw binary process, so divergence rates are preserved exactly. The many-to-one character of $\Fold_m$ on isolated windows therefore records only local aliasing once overlap is taken into account.

\begin{corollary}[Path-space equivalence and divergence-rate preservation]\label{cor:path-space-divergence}
Assume $m\ge 3$, let $\{\nu_u:u\in\mathcal U\}$ be a family of shift-invariant Borel probability measures on $\{0,1\}^{\mathbb Z}$, and write
\[
\eta_u:=(\Phi_m)_*\nu_u
\]
for their images on $Y_m$. Then $\Phi_m$ is a measurable isomorphism between the path-space experiments $(\{0,1\}^{\mathbb Z},\nu_u)$ and $(Y_m,\eta_u)$. In particular, for every $u,v\in\mathcal U$, the finite-block Kullback--Leibler divergences agree exactly after the block identification by $\Phi_{m,n}$, so the relative entropy rates satisfy
\[
h(\nu_u\|\nu_v)=h(\eta_u\|\eta_v)
\]
whenever these specific relative entropy rates exist. If $\nu_v=\nu_u^{\mathrm{rev}}$ is the time-reversal of $\nu_u$ and the entropy-production rate is defined by $h(\nu_u\|\nu_u^{\mathrm{rev}})$, then the same identity holds for entropy-production rates.
\end{corollary}

\begin{proof}
The path-space isomorphism is exactly Theorem~\ref{thm:two-sided-conjugacy}(i)--(ii). For every $n\ge m$, the block map $\Phi_{m,n}$ identifies the length-$(n+m-1)$ raw marginals of $\nu_u$ and $\nu_v$ with the length-$n$ stabilized marginals of $\eta_u$ and $\eta_v$. The finite-block Kullback--Leibler divergences therefore coincide term by term. Dividing by $n$ and letting $n\to\infty$ yields the equality of specific relative entropies whenever these limits exist. The entropy-production statement is the special case in which the reference law is the time reversal of the forward process and the rate is defined through the corresponding specific relative entropy.
\end{proof}

\begin{proposition}[Localization of irreversibility above threshold]\label{prop:irreversibility-localization}
Assume $m\ge 3$. Let $\nu$ be a shift-invariant Borel probability measure on $\{0,1\}^{\mathbb Z}$, let $\eta:=(\Phi_m)_*\nu$, and let $\mu_m$ and $\pi_m$ denote the corresponding length-$m$ marginals on $\{0,1\}^m$ and $X_m$. Then:
\begin{enumerate}[label=(\roman*),leftmargin=*,itemsep=2pt]
  \item the finite-window entropy ledger of Theorem~\ref{thm:kl-ledger},
  \[
  H(\mu_m)=H(\pi_m)+\EE_{\pi_m}[\log d_m(X)]-D_{\KL}(\mu_m\|\mu_m^e),
  \]
  describes the unresolved multiplicity of isolated windows;
  \item the path-space experiments $(\{0,1\}^{\mathbb Z},\nu)$ and $(Y_m,\eta)$ are measurably isomorphic, so specific relative entropy and, when defined through time reversal, entropy-production rate are the same on both sides.
\end{enumerate}
Consequently canonical admissibilization above threshold does not generate a new extensive source of irreversibility on path space.
\end{proposition}

\begin{proof}
Item~(i) is exactly Theorem~\ref{thm:kl-ledger}, which is formulated at fixed window length $m$. Item~(ii) is the single-process specialization of Corollary~\ref{cor:path-space-divergence}. Since the path-space laws are measurably isomorphic, any extensive information loss would have to appear in the specific relative entropy or entropy-production rate; item~(ii) excludes this. Therefore the degeneracy term from item~(i) is a local aliasing effect of isolated windows rather than a process-level source of irreversibility.
\end{proof}

\subsection{Exact transport of continuous potentials and finite-range energetics}\label{sec:measurement:renorm}

This subsection records the main dynamical and thermodynamic consequences of the finite-memory conjugacy. By transport we mean only passage through the finite-memory conjugacy. No renormalization-group flow, fixed point, or universality claim is intended. The first result applies to arbitrary continuous potentials; the finite-range formulas that follow give the explicit local form relevant for local observables and benchmark Hamiltonians.

\begin{theorem}[Transport of continuous potentials under the conjugacy]\label{thm:continuous-potentials}
Assume $m\ge 3$, let $\varphi:\{0,1\}^{\mathbb Z}\to\RR$ be continuous, and define
\[
\widetilde\varphi:=\varphi\circ \Phi_m^{-1}:Y_m\to\RR.
\]
Then:
\begin{enumerate}[label=(\roman*),leftmargin=*,itemsep=2pt]
  \item for every shift-invariant Borel probability measure $\mu$ on $\{0,1\}^{\mathbb Z}$,
  \[
  \int_{Y_m}\widetilde\varphi\,d(\Phi_{m*}\mu)=\int \varphi\,d\mu,
  \qquad
  h_{\Phi_{m*}\mu}(\sigma)=h_\mu(\sigma);
  \]
  \item the topological pressures agree:
  \[
  P_{Y_m}(\widetilde\varphi)=P_{\{0,1\}^{\mathbb Z}}(\varphi);
  \]
  \item $\mu$ is an equilibrium state for $\varphi$ if and only if $\Phi_{m*}\mu$ is an equilibrium state for $\widetilde\varphi$.
\end{enumerate}
\end{theorem}

\begin{proof}
Theorem~\ref{thm:two-sided-conjugacy}(i) states that $\Phi_m$ is a topological conjugacy onto $Y_m$, so $\Phi_m^{-1}$ is continuous and $\widetilde\varphi$ is again continuous. The integral identity is immediate from the definition of pushforward. The entropy identity is Proposition~\ref{prop:measure-entropy-preservation}. Therefore for every invariant measure $\mu$,
\[
h_{\Phi_{m*}\mu}(\sigma)+\int \widetilde\varphi\,d(\Phi_{m*}\mu)
=
h_\mu(\sigma)+\int \varphi\,d\mu.
\]
Taking the supremum over invariant measures and using the variational principle on both systems yields the pressure identity. The equilibrium-state correspondence is then immediate from equality of the variational functional under pushforward.
\end{proof}

\begin{corollary}[Parameterized phase diagrams are transported exactly]\label{cor:phase-diagram}
Assume $m\ge 3$, and let $\{\varphi_\lambda:\lambda\in\Lambda\}$ be any family of continuous potentials on the raw shift. Then for every parameter value $\lambda$,
\[
P_{Y_m}(\widetilde\varphi_\lambda)=P_{\{0,1\}^{\mathbb Z}}(\varphi_\lambda),
\qquad
\widetilde\varphi_\lambda:=\varphi_\lambda\circ\Phi_m^{-1}.
\]
Hence every phase boundary or nonanalytic point defined through the pressure function, and every equilibrium-state correspondence defined by the variational principle, is preserved pointwise under stabilized readout.
\end{corollary}

\begin{proof}
Apply Theorem~\ref{thm:continuous-potentials} to each parameter value separately.
\end{proof}

\begin{proposition}[Finite-range energetics under local inversion]\label{prop:finite-range-renorm}
Assume $m\ge 3$. Let $\phi:\{0,1\}^r\to\mathbb R$ be a finite-range interaction of range $r\ge 1$ on the raw shift. Define the induced interaction on $Y_m$ by
\[
\widetilde{\phi}(y_t,\dots,y_{t+r+m-2})
:=
\phi\bigl(\psi_m(y_t,\dots,y_{t+m-1}),\dots,\psi_m(y_{t+r-1},\dots,y_{t+r+m-2})\bigr).
\]
Then $\widetilde{\phi}$ has range $r+m-1$ on $Y_m$, and for every $a\in\{0,1\}^{\mathbb Z}$, $y:=\Phi_m(a)$, and $n\ge 1$,
\[
\sum_{t=0}^{n-1}\phi(a_t,\dots,a_{t+r-1})
=
\sum_{t=0}^{n-1}\widetilde{\phi}(y_t,\dots,y_{t+r+m-2}).
\]
\end{proposition}

\begin{proof}
For every $t\in\mathbb Z$, Theorem~\ref{thm:two-sided-conjugacy}(ii) gives
\[
a_t=\psi_m(y_t,\dots,y_{t+m-1}).
\]
Substituting this identity into $\phi(a_t,\dots,a_{t+r-1})$ yields the definition of $\widetilde{\phi}$ and proves the equality of the finite-volume energy sums term by term.
\end{proof}

\begin{corollary}[Thermodynamic transport above threshold]\label{cor:thermo-equivalence}
Assume $m\ge 3$. For every finite-range interaction $\phi$ on the raw full shift, the induced interaction $\widetilde{\phi}$ on $Y_m$ satisfies:
\begin{enumerate}[label=(\roman*),leftmargin=*,itemsep=2pt]
  \item finite-volume Hamiltonians agree exactly on matched bi-infinite configurations and therefore differ by at most an $O(1)$ boundary term under any standard finite-volume boundary convention;
  \item the corresponding partition functions agree up to boundary factors independent of volume, so pressure and free-energy density coincide;
  \item equilibrium states are transported bijectively by $\Phi_m$.
\end{enumerate}
In this restricted sense the stabilized readout transports finite-range interactions exactly above the three-window threshold.
\end{corollary}

\begin{proof}
Item~(i) is Proposition~\ref{prop:finite-range-renorm}. Item~(ii) follows because the energy sums agree exactly on matched configurations, and any discrepancy between finite-volume conventions comes only from the fixed-width boundary needed to evaluate the range-$(r+m-1)$ potential on $Y_m$. The pressure and free-energy-density statements are obtained by dividing the logarithms of the partition functions by volume and letting the volume tend to infinity. Item~(iii) is Proposition~\ref{prop:conjugacy-consequences}(iii) applied to the potential $\widetilde{\phi}$ transported through the conjugacy.
\end{proof}

\begin{corollary}[Occupancy and nearest-neighbour energy under stabilized readout]\label{cor:occupancy-energy}
Assume $m\ge 3$ and interpret $A_t\in\{0,1\}$ as a binary occupancy readout along the scan. Define the finite-volume observables
\[
N_n(A):=\sum_{t=0}^{n-1}A_t,
\qquad
C_n(A):=\sum_{t=0}^{n-1}A_tA_{t+1},
\]
and the nearest-neighbour energy
\[
\mathcal H_n(A):=-J\,C_n(A)-h\,N_n(A).
\]
Then for $Y=\Phi_m(A)$ there exist local observables on the stabilized process,
\[
\widetilde A_t(Y):=\psi_m(Y_t,\dots,Y_{t+m-1}),
\qquad
\widetilde C_t(Y):=\widetilde A_t(Y)\widetilde A_{t+1}(Y),
\]
such that
\[
N_n(A)=\sum_{t=0}^{n-1}\widetilde A_t(Y),
\qquad
C_n(A)=\sum_{t=0}^{n-1}\widetilde C_t(Y),
\qquad
\mathcal H_n(A)=\widetilde{\mathcal H}_n(Y)
\]
with
\[
\widetilde{\mathcal H}_n(Y):=-J\sum_{t=0}^{n-1}\widetilde C_t(Y)-h\sum_{t=0}^{n-1}\widetilde A_t(Y).
\]
Consequently, occupancy density, pair density, and the corresponding free-energy density are preserved exactly by stabilized readout.
\end{corollary}

\begin{proof}
The identities for $N_n$ and $C_n$ follow immediately from Theorem~\ref{thm:two-sided-conjugacy}(ii), which recovers each occupancy bit as a local observable of the stabilized trajectory. The energy identity is the special case of Proposition~\ref{prop:finite-range-renorm} with range $r=2$.
\end{proof}

\subsection{Secondary information-geometric consequences}\label{sec:measurement:info}

This subsection is logically downstream from the conjugacy and thermodynamic transport results. Above threshold the block experiments are related by a literal bijection, not merely by a sufficient statistic, so the same deterministic equivalence also preserves the block-level information geometry.

\begin{corollary}[Exact transport of block-level information geometry]\label{cor:info-geometry}
Fix $m\ge 3$ and $n\ge m$.
\begin{enumerate}[label=(\roman*),leftmargin=*,itemsep=2pt]
  \item For any convex function $f:[0,\infty)\to\RR$ with $f(1)=0$, let
  \[
  D_f(P\|Q):=\sum_a Q(a)\,f\!\left(\frac{P(a)}{Q(a)}\right)
  \]
  be the associated Csiszar $f$-divergence on the finite raw block space, with the usual conventions at $Q(a)=0$. Then for any probability laws $P,Q$ on $\{0,1\}^{n+m-1}$,
  \[
  D_f(P\|Q)=D_f\bigl((\Phi_{m,n})_*P\|(\Phi_{m,n})_*Q\bigr).
  \]
  In particular, the Hellinger integrals and the Bhattacharyya coefficient are preserved exactly.
  \item Let $\{P_\theta:\theta\in\Theta\}$ be a family of strictly positive mass functions on $\{0,1\}^{n+m-1}$ such that $p_\theta(a):=P_\theta(\{a\})$ is $C^1$ in $\theta$ for every raw block $a$. Let $Q_\theta:=(\Phi_{m,n})_*P_\theta$ and write
  \[
  q_\theta(y):=Q_\theta(\{y\}),
  \qquad y\in\mathcal Y_{m,n}.
  \]
  Then
  \[
  q_\theta(y)=p_\theta(\Psi_{m,n}(y)),
  \qquad
  \nabla_\theta \log q_\theta(y)=\nabla_\theta \log p_\theta(\Psi_{m,n}(y)).
  \]
  Consequently the Fisher information matrices of the raw and stabilized block experiments are identical:
  \[
  \mathcal I_Y(\theta)=\mathcal I_A(\theta).
  \]
\end{enumerate}
\end{corollary}

\begin{proof}
For (i), $\Phi_{m,n}$ is a bijection between the raw block space and $\mathcal Y_{m,n}$, with inverse $\Psi_{m,n}$. Writing the defining sum for the $f$-divergence on $\mathcal Y_{m,n}$ and changing variables by $a=\Psi_{m,n}(y)$ gives the identity immediately. The Hellinger and Bhattacharyya quantities are special cases obtained by choosing the corresponding functions of the likelihood ratio.

For (ii), the mass-function identity
\[
q_\theta(y)=p_\theta(\Psi_{m,n}(y))
\]
is exactly the pushforward relation under a bijection. Differentiating with respect to $\theta$ gives the score identity. The Fisher information matrix is the expectation of the score outer product, so the same change of variables under $\Psi_{m,n}$ shows that the raw and stabilized expectations coincide.
\end{proof}


